% Part:sets-functions-relations
% Chapter: size-of-sets
% Section: zig-zag

\documentclass[../../../include/open-logic-section]{subfiles}

\begin{document}

\olfileid{sfr}{siz}{zigzag}

\olsection{De zigzagmethode van Cantor}

\begin{explain}
We hebben al enkele ``makkelijke'' manieren gezien om een verzameling in een
lijst te zetten. Nu bekijken we een geval dat wat lastiger is. Neem de
verzameling paren van natuurlijke getallen\oliflabeldef{sfr:set:pai:sec}{, die
we in \olref[set][pai]{sec} zo hebben vastgelegd:}{, die als volgt is vastgelegd:}
\[
\Nat \times \Nat = \Setabs{\tuple{n,m}}{n,m \in \Nat}
\]
We kunnen deze geordende paren in een \emph{tabel} zetten:
\[
\begin{array}{ c | c | c | c | c | c}
& \mathbf 0 & \mathbf 1 & \mathbf 2 & \mathbf 3 & \dots \\
\hline
\mathbf 0 & \tuple{0,0} & \tuple{0,1} & \tuple{0,2} & \tuple{0,3} & \dots \\
\hline
\mathbf 1 & \tuple{1,0} & \tuple{1,1} & \tuple{1,2} & \tuple{1,3} & \dots \\
\hline
\mathbf 2 & \tuple{2,0} & \tuple{2,1} & \tuple{2,2} & \tuple{2,3} & \dots \\
\hline
\mathbf 3 & \tuple{3,0} & \tuple{3,1} & \tuple{3,2} & \tuple{3,3} & \dots \\
\hline
\vdots & \vdots & \vdots & \vdots & \vdots & \ddots\\
\end{array}
\]
Elk geordend paar uit $\Nat \times \Nat$ staat precies één keer in de tabel.
Het paar $\tuple{n,m}$ staat op rij $n$ en in kolom $m$. Maar hoe maken we van
alle vakken één lange lijst? De volgende tabel laat één manier zien. Er zijn
ook veel andere manieren:
\[
\begin{array}{ c | c | c | c | c | c | c}
& \mathbf 0 & \mathbf 1 & \mathbf 2 & \mathbf 3 & \mathbf 4 &\dots \\
\hline
\mathbf 0 & 0  & 1& 3 & 6& 10 &\ldots \\
\hline
\mathbf 1 &2 & 4& 7 & 11 & \dots &\ldots \\
\hline
\mathbf 2 & 5 & 8 & 12 & \ldots & \dots&\ldots \\
\hline
\mathbf 3 & 9 & 13 & \ldots & \ldots & \dots & \ldots \\
\hline
\mathbf 4 & 14 & \ldots & \ldots & \ldots & \dots & \ldots \\
\hline
\vdots & \vdots & \vdots & \vdots & \vdots&\ldots & \ddots\\
\end{array}
\]\noindent
Dit patroon heet \emph{de zigzagmethode van Cantor}. Zo zetten we alle paren uit
$\Nat \times \Nat$ in deze volgorde:
\[
\tuple{0,0}, \tuple{0,1}, \tuple{1,0}, \tuple{0,2}, \tuple{1,1},
\tuple{2,0}, \tuple{0,3}, \tuple{1,2}, \tuple{2,1}, \tuple{3,0}, \dots
\]
Daaruit volgt:
\end{explain}

\begin{prop}\ollabel{natsquaredenumerable}
De verzameling $\Nat \times \Nat$ is !!{enumerable}.
\end{prop}

\begin{proof}
Neem de functie $f \colon \Nat \to \Nat\times\Nat$. Die koppelt elk getal
$k \in \Nat$ aan het paar $\tuple{n,m} \in \Nat \times \Nat$ waarvoor $k$ in
Cantors zigzagtabel op rij $n$ en in kolom $m$ staat. 
\end{proof}

\begin{explain}
Deze aanpak laat zich ook mooi uitbreiden. Zo kunnen we alle geordende
drietallen van natuurlijke getallen in een lijst zetten. In zo'n drietal staan
drie getallen in een vaste volgorde:
\[
\Nat \times \Nat \times \Nat = \Setabs{\tuple{n,m,k}}{n,m,k \in \Nat}
\]
We bekijken $\Nat \times \Nat \times \Nat$ als het cartesische product van
$\Nat \times \Nat$ met $\Nat$. Dat betekent:
\[
\Nat^3 = (\Nat \times \Nat) \times \Nat =
\Setabs{\tuple{\tuple{n,m},k}}{n, m, k
  \in \Nat }
\]
Zo kunnen we ook $\Nat^3$ met een tabel in één lijst zetten. Langs de ene as
zetten we de lijst van de getallen in $\Nat$ en langs de andere de lijst van de
paren in $\Nat^2$:
\[
\begin{array}{ c | c | c | c | c | c}
& \mathbf 0 & \mathbf 1 & \mathbf 2 & \mathbf 3 & \dots \\
\hline
\mathbf{\tuple{0,0}} & \tuple{0,0,0} & \tuple{0,0,1} & \tuple{0,0,2} & \tuple{0,0,3} & \dots \\
\hline
\mathbf{\tuple{0,1}} & \tuple{0,1,0} & \tuple{0,1,1} & \tuple{0,1,2} & \tuple{0,1,3} & \dots \\
\hline
\mathbf{\tuple{1,0}} & \tuple{1,0,0} & \tuple{1,0,1} & \tuple{1,0,2} & \tuple{1,0,3} & \dots \\
\hline
\mathbf{\tuple{0,2}} & \tuple{0,2,0} & \tuple{0,2,1} & \tuple{0,2,2} & \tuple{0,2,3} & \dots\\
\hline
\vdots & \vdots & \vdots & \vdots & \vdots & \ddots \\
\end{array}
\]
Met een aanpak zoals de zigzagmethode van Cantor kunnen we dus ook alle
elementen van~$\Nat^3$ in een lijst zetten. We kunnen hiermee doorgaan en de
elementen van $\Nat^n$ in een lijst zetten, voor elk natuurlijk getal $n$.
Dus:
\end{explain}

\begin{prop}
$\Nat^n$ is !!{enumerable} voor elke $n \in \Nat$.
\end{prop}

\begin{prob}\label{sfr:siz:zigzag:prob:posint-n}
Laat zien dat $(\PosInt)^n$ !!{enumerable} is voor elke $n \in \Nat$.
\end{prob}

\begin{prob}\label{sfr:siz:zigzag:prob:posint-star}
  Laat zien dat $(\PosInt)^*$ !!{enumerable} is. Je mag
  \cref{sfr:siz:zigzag:prob:posint-n} gebruiken.
\end{prob}
\end{document}